\chapter{Assignment 1 --- WCET}

\section{Assignment}
After learning about scheduling and worst-case execution time (WCET) analysis, the first assignment focuses on developing a simple real-time application and analyzing its scheduling behavior and WCET. A passenger protection system is a system that has strong requirements for real time and fault tolerance. The following tasks should be completed: 
\begin{itemize}
    \item Implement a simplified version of the Airbag Deployment Algorithm described in \cite{mahmudNewDecisionMaking1995}. Using either C or C++, Don't use floating point numbers or OO features. 
    \item Analyze the WCET using IAR embedded workbench.
    \item Use BoundT to analyze the WCET. compare both results.
\end{itemize}

\section{Implementation}
The provided paper and lecture slides provide a starting point and outline how to implement the airbag deployment algorithm. 
It works by consistently calculating the current change in velocity as well as the length of the acceleration graph. Based in these values the system enters a STANDBY state and fires the airbag if the right conditions are met. 
Three thresholds are defined:  
\begin{itemize}
    \item $\Delta V_1$: The change in velocity must be greater than $\Delta V_1$ to enter the STANDBY state.
    \item $\Delta V_2$: If change in velocity exceeds this value the airbag is fired immediately. $\Delta V_2$ must be greater than $\Delta V_1$.
    \item $L$: The length of the acceleration graph must be greater than $L$ during standby state to fire the airbag.
\end{itemize}

Before implementation the algorithm was broken down into steps and states as follows: 

\subsection{Initialization} Sliding windows are initialized as arrays of size $N$ and filled with zeros. The system starts in the \emph{NORMAL} state. Everything after this is executed in an endless loop. Rudimentary timing is implemented using \texttt{usleep(1000000 / SAMPLES\_PER\_SECOND)}.

\begin{CCode}
    // Initialize arrays for storing time samples
    int t1[T_SIZE1] = {0}; // Array for first time series
    int t2[T_SIZE2] = {0}; // Array for second time series

    // Define thresholds for change in velocity
    int deltaV1 = 20; // Threshold for entering standby state
    int deltaV2 = 50; // Threshold for firing the airbag
    int LThreshold = 300; // Threshold for length of acceleration graph

    // Variables for current calculations
    int deltaV = 0; // Current change in velocity
    int L = 0; // Current length of acceleration graph

    // Variable to hold the current sample index
    int currentSample;

    // State of the airbag system
    AirbagState state = NORMAL; // Initial state is NORMAL
\end{CCode}

\subsection{Data Acquisition} Query new data, random values in this case, and update the sliding windows by shifting existing values and adding the new data. This step happens regardless of state. 

\begin{CCode}
    static void processSample(int t1[], int size1, int t2[], int size2, int *deltaV, int *L)
    {
    int currentSample = rand() % (10 + 1 + 30) - 30;
    printf("Current Sample: %d\n", currentSample);

    slideWindow(t1, size1, currentSample);
    slideWindow(t2, size2, currentSample);

    *deltaV = determineDeltaV(t1, size1);
    *L = determineL(t2, size2);

    printf("DeltaV: %d, L: %d\n", *deltaV, *L);
}

\end{CCode}

\subsection{Calculation} Calculate $\Delta V$ and $L$ within the scope of the sliding windows, using the formulas presented in the paper. 

\begin{CCode}
    static int determineDeltaV(int window[], int size)
    {
    double deltaV = 0;

    for (int i = 0; i < size; i++)
    {
        deltaV += window[i] * 1.0 / SAMPLES_PER_SECOND;
    }
    return round(deltaV);
}

static int determineL(int window[], int size)
    {

    double l = 0, dy = 0;
    double dx = 1.0 / SAMPLES_PER_SECOND;
    for (int i = 0; i < size - 1; i++)
    {
        dy = window[i + 1] - window[i];
        l += sqrt(pow(dy, 2) + pow(dx, 2));
    }
    return (int)l;
}
\end{CCode}

\subsection{State Transition} Possible states are \emph{NORMAL}, \emph{STANDBY} and \emph{FIRE}. If $\Delta V$ exceeds $\Delta V_1$ the system enters the standby state. If $\Delta V$ exceeds $\Delta V_2$ the system fires immediately. If the system is in STANDBY state and $L$ exceeds the threshold $L$, the system fires. If the system is in STANDBY state and $\Delta V$ drops below $\Delta V_1$ the system goes back to normal. For clarity these are implemented as switch cases.

Possible transitions are: 
\begin{itemize}
\item NORMAL $\rightarrow$ STANDBY $\rightarrow$ NORMAL 
\item NORMAL $\rightarrow$ STANDBY $\rightarrow$ FIRE
\end{itemize}

\begin{CCode}
        case NORMAL:
            processSample(t1, T_SIZE1, t2, T_SIZE2, &deltaV, &L);
            if (abs(deltaV) > deltaV1)
            {
                state = STANDBY;
            }
            break;

        case STANDBY:
            processSample(t1, T_SIZE1, t2, T_SIZE2, &deltaV, &L);
            if (abs(deltaV) > deltaV2 || L > LThreshold)
            {
                state = FIRE;
            }
            if (abs(deltaV) < deltaV1)
            {
                state = NORMAL;
            }
            break;

        case FIRE:
            exit(0);
            break;  

\end{CCode}

\subsection{Challenges} 
While the algorithm was simple to implement, trying to run it in the provided version of IAR Embedded Workbench caused some problems. Headers like time.h were not included per default and would have taken up too much space to compile in the free version. std:out was initially replaced with \texttt{printf} and later entirely omitted due to the same space concerns. Some rewrites were necessary to fix these problems.These problems were tackled as follows: 
\begin{itemize}
    \item unsupported headers: Removed time.h since only the WCET of the actual code execution matters for WCET, 
    \item Error e89: Caused by exceeding the 4KB object code limit of the kick start version. Remove the studio.h header to save on object code. This also means removing all \texttt{printf} statements, which did not matter for the WCET analysis anyway.
    \item Error e16: Occurred during building. Was fixed by selecting the \textit{ATmega328P} as the target processor. 
\end{itemize}
BoundT presented it's own set of challenges. It would, despite my best efforts, not compile on my machine so the pre-installed version in the proved VM was used instead. The 1024kb limit was no problem. The require .elf file was compiled from VS code using gcc, which failed due to missing STABS symbols. Using an older version of the compiler and including stabs finally lead to a usable result. 



\section{Evaluation}
After getting \textit{Wayland} and \textit{VMware Workstation} to stop fighting using the true and trust method of copying a fix from Reddit, all the analysis was done on the provided version of Windows XP. 

\begin{figure}[h]
    \centering
    \includegraphics[width=0.9\textwidth]{IAR_02.png}
    \caption{Profiling view from IAR Embedded Workbench Simulator. }
    \label{fig:IAR01}
\end{figure}

\subsection {WCET analysis from the IAR Simulator Results}

Debugging the code in IAR resulted in the following execution times for one loop after the initial setup (seen in \ref{fig:IAR01}) for the function implemented. \texttt{slideWindow} gets called twice per loop because the sliding windows are implemented as two separate arrays. The Accumulated Time for \texttt{main} shows that it takes 3998 cycles to execute one loop. At a clock speed of 16MHz this would take 0.25ms to execute, well under the 30ms requirement for the airbag system. Such a short time is expected because the sensor reading is not accounted for and the sliding windows are small with only a length of 20 and 30 respectively. Increasing both by 10 already increases the time to 5300 cycles or 0.33ms.



\subsection{WCET results from BoundT} 
After compiling the algorithm for the Atmega128 (the Atmega328P was not supported in BoundT) the analysis was run with the following parameters: 

\begin{CCode}
    boundt_avr -device=ATmega128 -elf_sym -max_steps 1024 -show times airbag_atmega128_stabs.elf main 
\end{CCode}

The resulting WCET in CPU cycles were: 
\begin{itemize}
    \item 2386 cycles for the \texttt{processSample} function, which coposes fetching the samples, sliding the windows and calculating $\Delta V$ and $L$.
    \item 3144 cycles for the \texttt{main} function, which includes the state transitions and the calls to \texttt{processSample}. 
\end{itemize}

These results are lower than the ones obtained from the IAR Simulator. This difference could come from the fact that the code was compiled for a different processor, or that the compilation in VS code performed some optimizations that the IAR compiler did not. 

\subsection{Discussion}

For both results it is important to keep in mind that the WCET from this example is not realistic since the code does not include the time it takes to read the sensors, which is a crucial part of the airbag system. The results are only for the processing of a small set of sample data and the state transitions, to satisfy the goal of this assignment which was to get familiar with the tools and the process of WCET analysis.
